\documentclass[11pt]{article}
\usepackage{ap-lang}
\title{Part 4: Style: Question Set}
\author{Paul Wang}
\date{November 2026}
\docpart{Part 4 \textperiodcentered\ Question Set}

\begin{document}
\maketitle

\begin{abstract}
Three exam-style passages with multiple-choice questions, followed by a second full rhetorical analysis essay (Question 2). Passage 1 is a personal essay about the disappearing night sky; Passage 2 is a satirical column in mock defense of meetings. Both reward close attention to diction, syntax, comparison, tone, and irony, and to what each choice does for the argument. Passage 3 is a writing set on a student's letter to the editor that needs revision for register, concision, and emphasis. All passages, writers, places, and figures are fictional and were written for this practice set. Allow about 35 minutes for the multiple choice and 40 minutes for the essay.
\end{abstract}

\scoringline{MCQ: \blank\ / \qtotal \qquad Essay: \blank\ / 6 \qquad Target: 19+ MCQ and 4+ essay \qquad Pace: 80 s per MCQ}

\section{Passage 1: The river of stars (reading)}

\begin{passage}{Passage 1}{The following passage is from a personal essay published in 2022 in a literary magazine.}
\pn{1} When I was nine, my father woke me at two in the morning, wrapped me in the wool blanket from the back of the truck, and carried me out past the barn to see the Milky Way. I had seen it before, I thought. I had not. What I had seen was a smudge. What he showed me that night was a river: thick, braided, so bright in places that it seemed to cast a shadow on the snow.

\pn{2} I went back to that field last winter. The barn is gone, and so is the dark. A distribution warehouse twelve miles east lights its parking lot all night, every night, for trucks that arrive twice an hour. Its glow sits on the horizon like a second, sickly dawn that never finishes rising. The river is a smudge again.

\pn{3} We have a word for the pollution of water and a word for the pollution of air, and we have laws for both. For the pollution of darkness we have a phrase, ``light pollution,'' that sounds faintly whimsical, like a problem for poets. It is not. Migrating birds, which steer by the stars, circle lit towers until they fall from exhaustion. Sea turtle hatchlings, which crawl toward the brightest horizon, crawl inland toward parking lots and die there. Our own bodies, which keep time by the coming of night, lose track of it.

\pn{4} The remedy is almost embarrassingly simple. Light that points down instead of up. Light that turns off when no one is there. Light that is warm rather than blue-white. None of this requires us to live in the dark; it requires us only to stop throwing light where no one needs it: into the sky, which has no use for it, and into the eyes of everyone who might, on some clear night, have looked up.

\pn{5} My father has been dead for six years. I would like my daughter to see what he showed me. I would like it not to require a four-hour drive.
\end{passage}

\begin{q}{2}{4.3 Comparisons}{``Smudge'' vs.\ ``river'' contrasts a dim impression with the overwhelming reality.}
In paragraph 1, the contrast between ``a smudge'' and ``a river'' primarily conveys
\ch{the cold of the winter night}{the difference between a faint impression of the Milky Way and its overwhelming presence under a truly dark sky}{the writer's confusion about astronomy as a child}{the writer's fear of the dark}
\end{q}

\begin{q}{3}{4.2 Syntax}{Two short sentences abruptly correct the child's assumption, marking a revelation.}
The sentences ``I had seen it before, I thought. I had not.'' (paragraph 1) mainly serve to
\ch{introduce the writer's father}{suggest the writer's memory is unreliable}{abruptly correct the child's assumption, marking the night as a revelation}{emphasize how often the writer had seen the Milky Way}
\end{q}

\begin{q}{3}{4.3 Comparisons}{The simile casts the glow as unnatural, unhealthy, and permanent.}
The simile ``like a second, sickly dawn that never finishes rising'' (paragraph 2) primarily characterizes the warehouse's glow as
\ch{beautiful but temporary}{a sign of economic progress}{unnatural, unhealthy, and permanently intrusive}{too faint to notice}
\end{q}

\begin{q}{1}{3.3 Structure}{``The river is a smudge again'' reverses the opening image to mark the loss.}
The sentence ``The river is a smudge again'' (paragraph 2) is effective chiefly because it
\ch{reverses the image from paragraph 1 to mark what has been lost}{shifts the essay's focus to the warehouse's business}{concedes that the opening memory was exaggerated}{introduces a new metaphor}
\end{q}

\begin{q}{4}{4.1 Diction}{She anticipates that the phrase invites readers not to take the issue seriously, then insists it is serious.}
In paragraph 3, the writer's remark that ``light pollution'' sounds ``faintly whimsical, like a problem for poets'' primarily serves to
\ch{suggest that poets are responsible for the problem}{propose that the phrase be replaced with a scientific term}{mock readers who care about the environment}{acknowledge that the term may lead readers to dismiss the issue before she insists on its seriousness}
\end{q}

\begin{q}{2}{4.2 Syntax}{Parallel ``which'' clauses give equal weight to harms across species, building a cumulative case.}
In paragraph 3, the parallel structure of the sentences about birds, hatchlings, and ``our own bodies'' primarily serves to
\ch{suggest that birds are harmed more than people are}{present a series of harms with equal weight, building a cumulative case that ends with the reader}{show that the writer is uncertain of the evidence}{introduce a list of possible solutions}
\end{q}

\begin{q}{3}{4.2 Syntax}{Fragments and repetition make the remedies sound short, concrete, and easy.}
In paragraph 4, the series of fragments beginning ``Light that\ldots'' primarily serves to
\ch{concede that the remedies would be expensive}{suggest that the remedies are too technical for most readers}{emphasize how simple and concrete the remedies are}{create a tone of despair}
\end{q}

\begin{q}{1}{4.4 Tone}{The turn to her father and daughter grounds the argument in personal stakes.}
The shift in the final paragraph primarily serves to
\ch{ground the public argument in the writer's personal loss and hope}{change the subject from light to family history}{introduce a counterargument to the essay's claim}{provide a statistic about driving distances}
\end{q}

\begin{q}{4}{4.4 Tone}{Mourning for what is lost, insistence that it can be recovered.}
The tone of the passage as a whole is best described as
\chh{playful and lighthearted}{bitter and despairing}{detached and technical}{elegiac and insistent}
\end{q}

\section{Passage 2: In defense of the meeting (reading)}

\begin{passage}{Passage 2}{The following passage is from a satirical column published in 2023 in a business magazine.}
\pn{1} It has become fashionable to complain about meetings. Critics say there are too many of them, that they run long, that they could have been an email. I would like to rise---slowly, as one does after ninety minutes in a conference chair---in their defense.

\pn{2} Consider what a meeting accomplishes. Twelve people, each of whom was in the middle of doing something, stop doing it. This alone is a great service; much of what we do in the middle of the afternoon is better left undone. They gather in a room whose thermostat no one controls and hear a summary of a document they were all sent and none of them read. In this way the document is, at last, read, at a cost of only twelve hours of combined salary.

\pn{3} The meeting also produces decisions, or, more precisely, it produces the next meeting, at which decisions will be made. This is prudent. A decision made too quickly might turn out to be correct, and then where would we be?

\pn{4} Some reformers propose that meetings have agendas, that they end on time, and that people who have nothing to contribute be allowed not to attend. I admire their idealism. But a meeting with an agenda is merely a list, and a meeting that ends on time is merely an interruption. Only the meeting without end achieves the true purpose of the institution, which is to ensure that everyone is too busy to notice that nothing is happening.

\pn{5} So the next time you are invited to a meeting about the meeting, accept. Bring a pen. You will not need it, but it will look as though you might.
\end{passage}

\begin{q}{1}{4.5 Punctuation}{The aside between dashes undercuts the defense with the critics' own complaint.}
In paragraph 1, the phrase set off by dashes (``slowly, as one does after ninety minutes in a conference chair'') primarily serves to
\ch{undercut the writer's stated defense with a detail that confirms the critics' complaint}{introduce the column's thesis}{shift the column to a serious tone}{establish the writer's expertise in office furniture}
\end{q}

\begin{q}{3}{4.4 Irony}{Praising people for stopping work is mock praise that criticizes.}
The statement that twelve people stopping their work is ``a great service'' (paragraph 2) is best described as
\chh{understatement}{an analogy}{verbal irony}{a concession}
\end{q}

\begin{q}{4}{4.1 Diction}{``Only'' ironically minimizes a large waste, drawing attention to it.}
In paragraph 2, the word ``only'' in ``at a cost of only twelve hours of combined salary'' primarily serves to
\ch{praise the company for paying its workers fairly}{suggest that the meeting should have been longer}{show that meetings are inexpensive}{ironically minimize a considerable waste, drawing attention to it}
\end{q}

\begin{q}{2}{4.2 Syntax}{The rhetorical question mocks an institutional preference for delay over results.}
The question ``A decision made too quickly might turn out to be correct, and then where would we be?'' (paragraph 3) primarily serves to
\ch{ask the reader to propose a solution}{mock institutions that prefer delay to results}{introduce statistics about decision-making}{express genuine concern about hasty decisions}
\end{q}

\begin{q}{1}{4.4 Irony}{The writer pretends to dismiss sensible reforms, implicitly endorsing them.}
In paragraph 4, the writer's treatment of the reformers' proposals is best described as
\ch{ironic dismissal that implicitly endorses the reforms}{a concession that the writer was wrong}{a sincere rejection of agendas and time limits}{an objective summary of competing views}
\end{q}

\begin{q}{2}{1.3 Purpose}{A mock defense is the vehicle for criticizing wasteful meetings.}
The writer's primary purpose in the column is to
\ch{argue that employees should refuse all meetings}{criticize wasteful meetings by pretending to defend them}{explain how to run a productive meeting}{defend the practice of holding long meetings}
\end{q}

\begin{q}{3}{4.4 Irony}{The pen is for appearances, which is the column's point about meetings.}
The final two sentences of the passage primarily serve to
\ch{concede that some meetings are useful}{suggest that the writer has changed his mind}{end with ironic advice that presents meetings as a performance of work rather than work}{offer practical advice for taking notes}
\end{q}

\begin{q}{4}{4.4 Tone}{Deadpan mock praise throughout: dry and satirical.}
The tone of the passage is best described as
\chh{earnest and admiring}{angry and accusatory}{anxious and uncertain}{dry and satirical}
\end{q}

\section{Passage 3: The Hollis Hardware mural (writing)}

\begin{passage}{Passage 3}{The following is a draft of a letter to the editor that a student is writing to a local newspaper. It may need revision.}
\pn{1} \sn{1} I am writing in response to last week's editorial calling the mural on the side of Hollis Hardware ``an eyesore.'' \sn{2} As one of the fourteen students who painted it, I would like to offer a different view.

\pn{2} \sn{3} The wall was not blank before we painted it. \sn{4} It was covered in graffiti that the store's owner, Mr.\ Hollis, had painted over at least nine times in five years. \sn{5} He asked our art class to paint something because, as he put it, ``nobody tags a wall that already says something.''

\pn{3} \sn{6} The mural shows the history of Fourth Street, from the streetcar line of the 1920s to the farmers market that fills the block on Saturdays today. \sn{7} Before we started, we first interviewed twenty-two longtime residents beforehand, and many of the faces in the mural are theirs. \sn{8} Mrs.\ Delgado, who ran the bakery on the corner for forty years, is the woman holding the bread.

\pn{4} \sn{9} The editorial said that the colors are ``garish.'' \sn{10} Okay, they are bright. \sn{11} We chose them on purpose, so that the mural could be seen from the bus stop across the street.

\pn{5} \sn{12} Since we finished in May, the wall has not been tagged once. \sn{13} People stop to find their neighbors in it. \sn{14} I think that is what a public wall is for.
\end{passage}

\begin{q}{1}{4.5 Punctuation}{A colon introduces what the wall was covered in, explaining sentence 3.}
The writer wants to combine sentences 3 and 4 in a way that emphasizes the explanation in sentence 4. Which choice best accomplishes this goal?
\ch{The wall was not blank before we painted it: it was covered in graffiti that the store's owner, Mr.\ Hollis, had painted over at least nine times in five years.}{The wall, not blank, before we painted it, was covered in graffiti, that the store's owner painted over.}{The wall was not blank before we painted it, it was covered in graffiti that the store's owner, Mr.\ Hollis, had painted over at least nine times in five years.}{The wall was not blank before we painted it and it was covered and Mr.\ Hollis had painted over the graffiti at least nine times in five years.}
\end{q}

\begin{q}{3}{2.2 Evidence}{The owner's own words show the mural was his remedy for graffiti.}
The writer is considering deleting the quotation in sentence 5 and ending the sentence after ``something.'' Should the writer make this change?
\ch{Yes, because quotations are not appropriate in a letter to the editor.}{Yes, because the quotation repeats information in sentence 4.}{No, because the owner's own words show that the mural was his solution to the graffiti, which supports the writer's view.}{No, because the quotation proves that the editorial's writer is wrong about the colors.}
\end{q}

\begin{q}{4}{4.5 Writing: style}{``Before we started,'' ``first,'' and ``beforehand'' say the same thing three times.}
Which version of the underlined portion of sentence 7 (reproduced below) is the most concise?

\emph{\underline{Before we started, we first interviewed twenty-two longtime residents beforehand},} and many of the faces in the mural are theirs.
\ch{Before we started, we first interviewed twenty-two longtime residents previously}{Before we started first, we interviewed twenty-two longtime residents}{Before we started, we initially first interviewed twenty-two longtime residents}{Before we started, we interviewed twenty-two longtime residents}
\end{q}

\begin{q}{2}{4.3 Imagery}{A specific visual detail lets readers who have not seen the mural picture it.}
The writer wants to add a sentence after sentence 8 that helps readers who have not seen the mural picture it. Which choice best accomplishes this goal?
\ch{The mural is very colorful and interesting.}{Behind her, a streetcar painted the deep green of the old Fourth Street line rattles toward the river.}{Murals are a popular form of public art.}{Many people have said they like the mural.}
\end{q}

\begin{q}{2}{4.5 Writing: style}{``Admittedly'' concedes the point in a formal register.}
The writer wants to revise sentence 10 so that it concedes the editorial's point while keeping the letter's formal tone. Which version best accomplishes this goal?
\ch{(as it is now)}{Admittedly, they are bright.}{Sure, whatever, they are bright.}{They are bright, so what?}
\end{q}

\begin{q}{1}{3.4 Transitions}{The new sentence signals the turn to the strongest evidence.}
The writer wants to add a sentence at the beginning of the last paragraph (before sentence 12) to signal a turn to her strongest evidence. Which choice best accomplishes this goal?
\ch{Most important, the mural has done exactly what Mr.\ Hollis hoped it would.}{Murals are popular in many cities.}{Also, there are some other things to consider.}{On the other hand, graffiti can also be a form of art.}
\end{q}

\begin{q}{1}{4.1 Diction}{``Find their neighbors in it'' conveys a personal, communal connection.}
The writer wants sentence 13 to emphasize the community's personal connection to the mural. Which choice best accomplishes this goal?
\ch{(as it is now)}{People stop and look at it.}{Many people see it every day.}{It is noticed by people.}
\end{q}

\begin{q}{4}{4.1 Diction}{Removing the hedge ``I think'' makes the closing claim confident and emphatic.}
The writer wants to end the letter with a more confident, emphatic statement. Which version of sentence 14 best accomplishes this goal?
\ch{(as it is now)}{Maybe that is kind of what a wall could be for, I guess.}{In conclusion, walls are for murals.}{That is what a public wall is for.}
\end{q}

\section{Free response}

\begin{frq}{Question 2 \textperiodcentered\ Rhetorical Analysis \textperiodcentered\ 6 points \textperiodcentered\ 40 minutes}
\small
\textbf{Directions.} The following column was published in 1958 by Harriet Lowe, the editor of the \emph{Carroll Valley Register}, in the final edition of the small-town weekly newspaper before it closed. (Lowe and her newspaper are fictional.) Read the passage carefully.

\textbf{Prompt.} Write an essay that analyzes the rhetorical choices Lowe makes to convey to her readers the value of a local newspaper.

In your response you should do the following:
\begin{itemize}[leftmargin=1.5em,itemsep=1pt]
\item Respond to the prompt with a thesis that analyzes the writer's rhetorical choices.
\item Select and use evidence to support your line of reasoning.
\item Explain how the evidence supports your line of reasoning.
\item Demonstrate an understanding of the rhetorical situation.
\item Use appropriate grammar and punctuation in communicating your argument.
\end{itemize}
\end{frq}

\begin{passage}{Question 2 passage}{Harriet Lowe, final column, \emph{Carroll Valley Register}, 1958.}
\pn{1} This is the last edition of the \emph{Carroll Valley Register}. For ninety-one years this paper has come to your porch on Thursday, and after today it will not, and I have spent a month trying to find a way to say that which does not sound like an obituary. I have failed. So let it be one.

\pn{2} The \emph{Register} was born in 1867 in the back room of a harness shop, and it died this year of the usual causes: rising costs, falling advertising, and a reader who could get the news from the city faster and for less. I do not blame that reader. I was often that reader myself.

\pn{3} But I want to say, before the presses go quiet, what the city papers will not do. They will not tell you that the Oberg boy made the honor roll, or that the creamery is hiring, or that the bridge on County Road 9 will be closed on Tuesday. They will not print the minutes of the school board, all of them, including the dull parts, which are the parts where your money is spent. They will not send anyone to sit through the zoning hearing on a February night. Nobody from the city has ever sat through a zoning hearing in Carroll Valley. I have sat through two hundred.

\pn{4} A town without a newspaper is not a town without news. The news happens anyway. It is only that no one writes it down, and what is not written down is soon remembered wrong, and what is remembered wrong is soon argued about, and what is argued about without a record is decided by whoever argues loudest.

\pn{5} I have kept every edition. They are in the basement of the library, in ninety-one bound volumes, and anyone may read them. You will find in them your grandparents' wedding, the fire that took the Methodist church, the price of hogs in 1931, and the long, foolish, wonderful argument about whether to pave Main Street. You will find, in other words, yourselves.

\pn{6} Thank you for reading. Somebody should keep doing it.
\end{passage}

\makekey

\subsection*{Essay scoring guide (6 points)}
\small
The rows follow the published AP English Language rubric (thesis 1, evidence and commentary 4, sophistication 1). The ``earns credit'' column shows what a top response does with this passage.
\begin{compare}{Row}{Earns credit}{Common miss}
A. Thesis (1) & A defensible thesis that analyzes Lowe's choices in relation to her purpose, e.g.\ ``Writing the paper's obituary, Lowe refuses to blame her readers and instead catalogs the small, specific work only a local paper does, so that by the end its loss feels like the loss of the town's memory of itself, a loss her readers can still prevent.'' & ``Lowe uses imagery and repetition to show newspapers are important''; summarizing the column. \\
B. Evidence and commentary (4) & Specific evidence with consistent commentary, e.g.\ the obituary frame (paragraphs 1--2: ``born'' and ``died'') gives the paper a life her readers have shared; refusing to blame the reader (``I was often that reader myself'') keeps a grieving audience from becoming defensive; the anaphora ``They will not\ldots'' and the concrete local details (the Oberg boy, County Road 9) make the paper's value tangible; the short ``I have sat through two hundred'' contrasts her endurance with the city's absence; the polysyndeton chain in paragraph 4 enacts how memory decays into conflict without a record; ``You will find\ldots yourselves'' turns the archive into the community's identity; the final line passes responsibility to readers. 4: all claims supported and consistently explained. 3: some explained. 2: evidence with little reasoning. 1: summary. & Naming the devices (anaphora, polysyndeton) without explaining their effect; treating the column as an attack on city papers or on readers. \\
C. Sophistication (1) & Explains a tension Lowe manages: mourning the paper without guilt-tripping the readers who let it die, or how the elegiac tone carries an argument about civic accountability (the school board minutes, ``where your money is spent''); or traces how the column moves from personal grief to public warning to a charge to the readers; or sustains a vivid, controlled style. & One sentence calling the column ``nostalgic''; claims about all journalism today unrelated to the passage. \\
\end{compare}

\end{document}
